\chapter{绪论}

\section{课题研究的重要意义}

本课题来源于2010年中国气象局公益性气象行业专项“冰雪自动化观测技术研究”（项目编号：GYHY201006045）。

我国是雨雪、冰冻灾害性天气多发的国家，具有影响范围广、程度深、发生频次高的特点。2008年初，我国南方地区连续遭受低温雨雪冰冻极端天气过程袭击，导致高速封路、机场关闭、高压线路破坏，对电力、通信、交通和人民生命财产带来极大危害。2009年末至2010年3月初，我国北方十多个省市地区遭受多次雪灾，给交通运输造成严重影响，对人们的生产生活特别是农林畜牧业生产造成巨大损失。

………………

………………

………………

………………

\section{雪深自动测量的研究现状}

\subsection{几种常用的雪深测量方法}

首先，介绍传统的人工测量雪深的方法，指出此方法的优点以及不足。然后介绍由此而产生的雪深自动测量的方法，主要都有哪些测量方法，简要介绍各种方法的特点\cite{ref-wetland,ref-water}。

………………

………………

光扫描法

投光器和受光器彼此保持一定距离，口朝下安装在距离地面的某一高度上。受光器的方向固定，投光器则常对雪面左右扫描，投光器的光进入受光器的视野时，根据投光器的角度求出积雪深度。

双杆法

将方向相对的投光器和受光器装在有数米间距的直立的两根杆子上，使其可沿杆同时上下滑动。当投光器、受光器位于雪面下方时，来自投光器的光受到阻挡，当它们一超出雪面，受光器就接收到投光器的光。根据这时的高度求出积雪深度。

………………

………………

表格格式：

\begin{table}[htbp]
  \centering
  \caption{“正确”的雪深瞬时值的判断条件\cite{ref-standard,ref-lantian}}
  \begin{tabular}{ccccc}
    \toprule
    气象变量 & 下限 & 上限 & 存疑的变化速率 & 错误的变化速率 \\
    \midrule
    雪深 & 0 mm & 1500 mm & 20--60 mm & $>60$ mm \\
    \bottomrule
  \end{tabular}
\end{table}

\begin{table}[htbp]
  \centering
  \caption{试验测试用具及性能指标}
  \begin{tabular}{cccc}
    \toprule
    序号 & 名称 & 型号 & 主要性能指标 \\
    \midrule
    1 & 传感器 & A型 & 分辨率1 mm \\
    2 & 采集器 & B型 & 采样间隔1 min \\
    \bottomrule
  \end{tabular}
\end{table}

图片格式（本图为例样，请根据论文情况据实画图）：

\begin{figure}[htbp]
  \centering
  \fbox{\rule{0pt}{6cm}\rule{8cm}{0pt}}
  \caption{试验装置图}
\end{figure}

公式格式：

\begin{equation}
  E = mc^2
\end{equation}
